\section{Notation}

Im Folgenden greifen wir die von Redmond et al. definierte Notation ihres formalen Modells auf und überführen diese in die Form, die auch im Rest dieser Arbeit verwendet wird, damit dem Leser ein einheitliches Notationskonzept zur Verfügung steht.

\subsection{Komponenten}
Komponenten erlauben die Speicherung und Modifikation domänenspezifischer Werte.\\

\noindent
Sei $K$ ein Komponentenschema.\\

\noindent
Stehe weiter mit $I$ eine Indexmenge von Bezeichnern (\textit{Labels}) zur Verfügung
Wir fassen damit $K_i$ als Familie auf, die über $i \in I$ indexiert ist.\\

\noindent
Bezeichne $k_i \in K_i$ einen beliebigen Term des Typs $K_i$.\footnote{``Each $K_i$ is a type in the sub-language, \ldots ``}
Redmond et al. definieren hierzu -- informell -- eine Abbildung:\footnote{
``The component types come to us in a mapping from component label $i \in I$ to component type $K_i$, \ldots ``
}

\[
    k: I \rightarrow K_i
\]

\noindent
wobei $K_i$ durch das Label $i \in I$ determiniert ist.
Wie in der zitierten Arbeit nutzen auch wir im Folgenden die Indexnotation $k_i$ anstelle $k(i)$:

\[
    k_i \in K_i
\]

\noindent
Hierdurch wird mit $k_i$ ein beliebiger, isolierter, zustandsloser Wert aus dem Wertebereich der durch den Komponententypen repräsentierten Domäne $K_i$ dargestellt.
Entsprechend kann ein beliebiger Term $k_i$ formal mit mehreren Entitäten assoziiert sein, ohne, dass es eines zusätzlichen Index bedarf.\footnote{
    In der Praxis wirkt sich dies als Kopie der entsprechenden Werte aus (Redmond et al.: ``Component values``).
}\\

\subsubsection*{Beispiel}
Redmond et al. notieren in ihrer Arbeit beispielhaft ein Komponentenschema $K$, das folgende Paare $i \mapsto K_i$ beinhaltet:
\[
    K \coloneqq \{ i_1 \mapsto K_{i_1}, i_2 \mapsto K_{i_2} \}
\]

\noindent
mit
\[
    \begin{align}
        i_1 \coloneqq \text{Pos}\\
        i_2 \coloneqq \text{Vel}\\
        K_{i_1}, K_{i_2} \coloneq \mathbb{Z}
    \end{align}
\]

\noindent
Für einen Komponententyp $K_\texttt{Position} \subseteq \mathbb{R}^2$ können wir also schreiben:

\[
    k_{\texttt{Position}} = (1.0, 1.0) \in K_\texttt{Position}
\]



\begin{tcolorbox}[title={Komponenten als Informationsträger für Referenzen}]
Für Anwendungsfälle, in denen eine Komponente eine Referenz auf eine Objektinstanz $t$ vom Typ $T$ verwaltet, kann in diesem Kontext das Label

\[
    \text{Ref(T)} \in I
\]

\noindent
eingeführt werden.\\

\noindent
Hierzu sei

\[
    \mathcal{A} \colon \text{Inst}(T) \rightarrow \mathbb{N}
\]

\noindent
eine Abbildung, die eine Instanz eines Typs $T$ auf einen numerischen Wert abbildet -- den Speicherort der Instanz im \textit{Adressraum}.\\

\noindent
Ein Term $k_{\text{Ref(T)}}$ speichert dann diesen numerischen Wert für $t \in \text{Inst}(T)$

\[
    k_{\text{Ref(T)}} = \mathcal{A}(t) \in K_{\text{Ref(T)}} = \mathbb{N}
\]


\noindent
Dadurch ist die Frage, wie mit diesem Komponentenwert im ECS-Programm umgegangen wird, zunächst semantischer Natur, und unabhängig von der Tatsache, dass der Term tatsächlich einen Speicherort repräsentiert.
\end{tcolorbox}

\subsection{Entitäten}

Sei mit $E$ ein Typ gegeben, der eine Wertemenge\footnote{
    Beispielswiese \texttt{uint32\_t} in \texttt{C++}.
} zur Identifikation zusammengehöriger Komponentenbündel ermöglicht.\\

\noindent
Wenn der Term $e$ einen bestimmten Wert des Typs $E$ repräsentiert, dann ist
\[
    \mathcal{P} \coloneqq E \times \bigcup_{i \in I} K_i = \{(e, k_i) \mid e \in E, i \in I, k_i \in K_i\}
\]
die Menge der möglichen Entität-Komponenten-Paare.\\

\bigskip
Eine bestimmte, eindeutige Ausprägung einer Entität vom Typ $E$, die wir mit $\mathcal{E}_m, m \in \mathbb{N}$ notieren, und die damit verbundene mögliche Menge von Entität-Komponenten-Paaren kann mithilfe folgender partiellen Abbildung hergeleitet werden:
\[
    \mathcal{E}_m :  I \rightharpoonup \bigcup_{i \in I} K_i, \quad i \mapsto k_i \in K_i, m \in \mathbb{N}
\]
Hierbei gilt mit der Definition der partiellen Abbildung:
\[
    \mathcal{E}_m(i) \neq \bot \Leftrightarrow \exists{!} k_i \in K_i : \mathcal{E}_m(i) = k_i
\]

\noindent
Hiermit lässt sich die \textit{assoziierte Komponentenmenge} definieren:

\[
    \mathcal{E}^K_m \coloneqq \{k_i \mid k_i = \mathcal{E}_m(i), k_i \in K_i, i \in I \}
\]

\noindent
\textbf{Beispiel}:

\[
    \begin{alignat}{3}
    \mathcal{E}^K_{1} &= \{\mathcal{E}_{1}(i_1), \mathcal{E}_{1}(i_2), \ldots, \mathcal{E}_{1}(i_n)\}\\
            &= \{k_{i_1}, k_{i_2}, \ldots, k_{i_n}\}
    \end{alignat}
\]

\vspace{5mm}
\noindent
Mit der Definition von $\mathcal{E}_m$ gilt für verschiedene $m, n$ trivialerweise, dass zwei verschiedene Entitäten dieselbe Menge von Termen besitzen können -- damit lässt sich die Eindeutigkeit einer Entität nicht über ihre entsprechende Komponentenmenge feststellen.\\

Aus diesem Grund wird bei Redmond et al. diese Darstellung nicht verwendet.
Ganz im Sinne der Datenorientierung ist die Zuordnung der Entitäten zu Komponenten über Zustände definiert, die eine Menge von Abbildungen der Komponenten-Typen auf eine Menge von partiellen Abbildungen $e_m \rightarrow k_i$ enthält (siehe Abschnitt~\ref{sec:zustand}).
Dies erleichtert auch die formale Herleitung der Identität von Entitäten: Somit exitiert in dem späteren Modell immer nur genau eine eindeutig identifizierbare Entität-Komponentent-Konstellation.
\bigskip

\noindent
Redmond et al. binden deshalb die Eindeutigkeit einer Entität an den Term $e \in E$:

\[
    E \coloneqq \{e_1, e_2, \ldots\}
\]

\noindent
so dass alle Tupel in $\mathcal{P}$, die den Term $e_m$ enthalten, über die Familie

\[
    \mathcal{P}_{e_m} \coloneqq \{(e_m, k_i) \mid e_m \in E, \mathcal{E}_m(i) = k_i \in K_i\}, m \in \mathbb{N}
\]

\noindent
definiert werden kann.

\bigskip
\begin{theorem}\label{th:entity_id1}
    Gegeben sei $\mathcal{P}_{e_m}$ und $\mathcal{P}_{e_n}$, $e_n, e_m \in E$.
    Dann gilt:
    \[
         e_m \neq e_n \Rightarrow \mathcal{P}_{e_m} \neq \mathcal{P}_{e_n}
    \]
\end{theorem}
\bigskip

\begin{proof}
    Seien beliebige $x = (e_m, k_{i}) \in P_{e_m}, y = (e_n, k_j) \in P_{e_n}, k_i \in K_i, j \in K_j, i, j \in I$ gegeben.\\

    \noindent
    Mit der Voraussetzung $e_m \neq e_n$ gilt
    \[
        e_m \neq e_n \land (k_i = k_j \lor k_i \neq k_j)
    \]

    \begin{itemize}
        \item \textbf{Fall 1} - $e_m \neq e_n \land k_i = k_j$: Es ist $x(1) \neq y(1) \land x(2) = y(2)$ und damit $x \neq y$.
        \item \textbf{Fall 2} - $e_m \neq e_n \land k_i \neq k_j$: Es ist $x(1) \neq y(1) \land x(2) \neq y(2)$ und damit $x \neq y$.
    \end{itemize}

    \noindent
    In beiden Fällen ist  $\mathcal{P}_{e_n} \neq \mathcal{P}_{e_m}$ erfüllt.
\end{proof}
\bigskip

\noindent
Theorem~\ref{th:entity_id2} folgt direkt als Kontraposition zu Theorem~\ref{th:entity_id1}:\footnote{
    $A \Rightarrow B \equiv \neg B \Rightarrow \neg A$.
    Wir geben den Beweis trotzdem formal an, um die Bindung der Komponenten-Terme an den Entitätsterm zu verdeutlichen.
}

\begin{theorem}\label{th:entity_id2}
    Gegeben sei $\mathcal{P}_{e_m}$ und $\mathcal{P}_{e_n}$, $e_n, e_m \in E$.
    Dann gilt:
    \[
        \mathcal{P}_{e_m} = \mathcal{P}_{e_n} \Rightarrow e_m = e_n
    \]
\end{theorem}
\bigskip

\begin{proof}
    Seien beliebige $x = (e_m, k_{i}) \in P_{e_m}, y = (e_n, k_j) \in P_{e_n}, k_i \in K_i, j \in K_j, i, j \in I$ gegeben.\\

    \noindent
    Mit der Voraussetzung $\mathcal{P}_{e_m} = \mathcal{P}_{e_n}$ gilt
    \[
        \forall x \in \mathcal{P}_{e_m} : y \in \mathcal{P}_{e_n}
    \]

    \noindent
    Damit folgt direkt
    \[
        \begin{alignat}{3}
            x = y &\Leftrightarrow (e_m, k_i) = (e_n, k_j) \\
                  &\Leftrightarrow e_m = e_n \land k_i = k_j
        \end{alignat}
    \]
\end{proof}
\bigskip

\noindent
Wir haben damit eine triviale, aber wichtige Eigenschaft gezeigt: Die Identität einer Entität-Komponenten-Konstellation kann erst mit $e \in E$ definiert werden.
Dies wird später relevant, wenn es darum geht, den nebenläufigen Zugriff auf Entitäten und ihre Komponenten zu formalisieren.\\

Eine weitere wichtige Abgrenzung nehmen Redmond et al. bzgl. der assoziierten Komponententypen vor:

\[
    \forall (e_m, k_i), (e_m, k_j) \in \mathcal{P}_{e_m} : k_i \neq k_j \Rightarrow i \neq j
\]

\noindent
Eine Entität ist also mit höchstens einem Term $k_i$ pro Komponententyp $K_i$ assoziiert.
Davon unberührt bleiben die Werte der Terme: So kann auch weiterhin $k_i = k_j, i \neq j$ gelten.

\subsection{Zustand}
Den Begriff \textit{Zustand} (\textit{state}) führen Redmond et {al.} ein, um die zuvor erwähnte Eindeutigkeit von Entität-Komponenten-Konstellationen zu formalisieren.
Hierzu gilt wie bislang festgehalten, dass eine Entität aus Assoziationen zu unterschiedlichen, niemals paarweise gleichen Komponententypen besteht.\\

Um den Zustand eines \textit{Core ECS} Programms zu definieren, führen die Autoren $c$ ein: Eine über $i \in I$ indexierte Famielle von partiellen Abbildungen, die eine Entität mit einem Komponententerm vom Typ $K_i$ assoziieren:

\[
    c \coloneqq \{c_i : E \rightharpoonup K_i\}_{i \in I}
\]

\noindent
So lässt sich für eine Entität $e_m$ der derzeitige Wert des Komponententerms $k_i$ notieren über

\[
  c_i(e_m) = k_i
\]

\noindent
was äquivalent zu der oben hilfsweise eingeführten Definition von $\mathcal{E}_m$ ist:

\[
    c_i(e_m) = k_i = \mathcal{E}_m(i)
\]

\noindent
Wir behalten im weiteren die Notation von Redmond et {al.} bei, da sie im Sinne der Datenorientierung von einem Typ ausgehend den Wert für eine bestimmte Entität, die mit diesem Typ assoziert ist, abfragt\footnote{
    und gegebenefalls $\bot$ liefert, falls diese Entität mit dem Typ nicht assoziert ist
}.\\
Die Notation $\mathcal{E}_{m}$ geht -- als entitätszentrische Darstellung im gegensatz zu der komponentenzentrierten Darstellung von Redmond et al. -- zunächst einmal von der Existenz einer Entität aus, was im Sinne des Entity-Component-Systems nicht idiomatisch ist.
Außerdem sieht die Definition von \textit{Core ECS} gerade vor, dass ein Programm parametrisiert wird mit dem Typ $E$ und dem Komponentenschema $K$ - somit sind alle $i \in I$ bekannt, wobei während der Programmausführung weitere $e_m$ hinzukommen, und damit auch $E_m$ erweitert werden müsset.


\subsection{Systeme}